\documentclass[a4paper,12pt]{article}

\usepackage[utf8]{inputenc}
\usepackage[portuguese]{babel}
\usepackage{fancyhdr}
\usepackage{amsmath}
\usepackage{amsthm}
\usepackage{amssymb}
\usepackage[portuguese,algoruled,lined,linesnumbered]{algorithm2e}

% ---- PREENCHA ----
\newcommand{\Lista}{2}
\newcommand{\Nome}{Leonardo Heidi Almeida Murakami}
\newcommand{\NUSP}{11260186}
% -----------------------------

\pagestyle{fancy}
\fancyhf{}
\fancyhead[L]{\Nome}
\fancyhead[R]{NUSP \NUSP}
\fancyfoot[L]{Lista \Lista}
\fancyfoot[C]{MAC0338}

\begin{document}

\section*{Exercício 1}
\begin{algorithm}[H]
    \SetKwProg{Def}{def}{:}{}
    \SetAlgoLined
    \caption{CountRange}
    
    \KwData{lista $X[1..n]$ inteiros no intervalo 1 a k}
    \KwResult{numero $n$ de inteiros em um intervalo $[a..b]$}
    
    \Def{preprocess(X, n, k)}{
        $C \leftarrow \text{array [1..k]}$\;
        $P \leftarrow \text{array [1..k]}$\;
        \For{$i \leftarrow 1$ \KwTo k}{
            $C[i] \leftarrow 0$\;
        }
        \For{$i \leftarrow 1$ \KwTo n}{
            $val \leftarrow X[i]$\;
            $C[val] \leftarrow C[val] + 1$\;
        }
        $prefix \leftarrow 0$\;
        \For{$i \leftarrow 1$ \KwTo k}{
            $P[i] \leftarrow prefix$\;
            $prefix \leftarrow prefix + C[i]$\;
        }
        \Return{P}\;
    }
    \Def{countRange(P, a, b)}{
        \If{$a \ge b$}{
            \Return{0}\;
        }
        \If{$a == 1$}{
            \Return{$P[b]$}\;
        }
        \Return{$P[b]-P[a-1]$}\;
    }
\end{algorithm}
\section*{Exercício 2}
\begin{algorithm}[H]
\SetKwProg{Def}{def}{:}{}
\SetAlgoLined
\caption{sort}

\KwData{lista $X[1..(n^2-1)]$ nao ordenada}
\KwResult{lista $X[1..(n^2-1)]$ ordenada}
\Def{sort(X, n)}{
    \For{$i \leftarrow 1$ \KwTo $n$}{

    }
}
\end{algorithm}
\section*{Exercício 2}
\begin{algorithm}[H]
\SetKwProg{Def}{def}{:}{}
\SetAlgoLined
\caption{sort}

\KwData{lista $X[1..(n^2-1)]$ nao ordenada}
\KwResult{lista $X[1..(n^2-1)]$ ordenada}
\Def{sort(X, n)}{
    \For{$i \leftarrow 1$ \KwTo $n$}{

    }
}
\end{algorithm}
\section*{Exercício 6}
\begin{enumerate}
    \item Achamos o max do array, que é 6
    \item Inicializamos o array de contagem C[0..6]
    \item Fazemos o array de contagem passando uma vez por cada elemento de A e somando o valor de sua propria posição no array de contagem
    \item C[0..6] = [2, 2, 2, 2, 1, 0, 2]
    \item Transformamos o array de contagem em um array prefixados, acumulando a soma
    \item P[0..6] = [2, 4, 6, 8, 9, 9, 11]
    \item De tras para frente vemos quantos numeros menores que cada posicao existem
    \item O[1..11] = [_, _, _, _, _, 2, _, _, _, _, _]
    \item O[1..11] = [_, _, _, _, _, 2, _, 3, _, _, _]
    \item etc..
\end{enumerate}

\section*{Exercício 7}

\end{document}